% mini_expose_v3_politik.tex — Kurzfassung (max. 2 A4-Seiten)
% Kernaussage: Was will ich machen und warum.

\documentclass[11pt,a4paper]{scrartcl}
\usepackage[ngerman]{babel}
\usepackage[utf8]{inputenc}
\usepackage[T1]{fontenc}
\usepackage{lmodern}
\usepackage{csquotes}
\usepackage{microtype}
\usepackage{enumitem}
\usepackage{xcolor}
\usepackage{hyperref}
\usepackage[top=1.8cm,bottom=1.8cm,left=2.2cm,right=2.2cm]{geometry}

\hypersetup{colorlinks=true,linkcolor=black,urlcolor=blue!70!black}

\setlength{\parindent}{0pt}
\setlength{\parskip}{0.45em}
\setlist[itemize]{leftmargin=1.3em,topsep=0.2em,itemsep=0.15em}
\setlist[enumerate]{leftmargin=1.5em,topsep=0.2em,itemsep=0.15em}

% kompaktere Überschriften
\usepackage{titlesec}
\titlespacing*{\section}{0pt}{0.7em}{0.25em}
\titleformat{\section}{\normalfont\large\bfseries}{\thesection.}{0.5em}{}

\newcommand{\code}[1]{\texttt{#1}}

\begin{document}

\begin{center}
  {\Large\bfseries Mini-Exposé: Transparenz politischer Entscheidungsprozesse}\\[0.3em]
  {\small Nachvollziehbarkeit von Gesetzgebung als inspizierbarer Prozess und ihre Wirkung
  auf Vertrauen und demokratische Kontrolle \\
  Themenfeld: \enquote{Mehr Legitimität durch offene Verfahren? Transparenz und
  Beteiligung im politischen Willensbildungsprozess}}
\end{center}

\vspace{0.2em}
\hrule
\vspace{0.4em}

% ---------------------------------------------------------------------------
\section{Warum: das Problem}

In modernen Demokratien entscheidet die \textbf{Nachvollziehbarkeit politischer Verfahren}
über Vertrauen oder Entfremdung der Bürger. Zwei Probleme bleiben offen:

\begin{itemize}
  \item \textbf{Intransparent.} Zwischen Gesetzentwurf und Verabschiedung liegen zahllose
    Ausschusssitzungen, Anhörungen und Lobbykontakte, die für Außenstehende kaum sichtbar
    sind. Wer welche Formulierung eingebracht hat und \emph{warum} eine Regelung so und nicht
    anders lautet, lässt sich nachträglich selten rekonstruieren; der politische Prozess ist
    eine Black Box.
  \item \textbf{Wirkungslos beteiligend.} Formale Beteiligungsformate (Bürgerdialoge,
    Konsultationen) setzen oft erst \emph{nach} den entscheidenden Weichenstellungen an. Ob
    Bürger die Entstehung einer Regelung \textbf{vor} der Abstimmung bewerten und beeinflussen
    können und ob das die wahrgenommene Legitimität erhöht, ist bislang kaum systematisch
    untersucht, obwohl Bürger diese frühe Kontrolle nachweislich einfordern.
\end{itemize}

% ---------------------------------------------------------------------------
\section{Was wird umgesetzt:}

Ein digitales Verfahren, das den Weg eines Gesetzes nicht nur automatisch dokumentiert,
sondern die einzelnen Entscheidungsschritte vor der finalen Abstimmung inspizier- und
kommentierbar macht. Ein Gesetzgebungsprozess wird als nachvollziehbarer Entscheidungsgraph
(Akteure, Änderungsanträge, Begründungen, Abstimmungen) repräsentiert. Es gibt drei Bausteine.

\begin{enumerate}
  \item \textbf{Prozess-Rekonstruktion.} Verknüpfung von Entwurfsversionen, Ausschussprotokollen
    und Abstimmungsergebnissen über Zeit; jede Textänderung wird ihrem Anlass (Anhörung,
    Antrag, Stellungnahme) zugeordnet.
  \item \textbf{Entscheidungspfad.} Pro Regelung eine nachvollziehbare Begründung: welcher
    Akteur welche Formulierung einbrachte, mit welcher Begründung und welchem Abstimmungsverhalten.
  \item \textbf{Inspektions- und Beteiligungs-Interface.} Eine öffentliche Ansicht zeigt den
    geplanten Entscheidungspfad samt Begründungsspur. Bürger sehen \emph{warum} eine Regelung
    ihre Form annimmt, können einzelne Schritte kommentieren, Alternativen vorschlagen und den
    Stand verfolgen: Gesetzgebung als offener Vorschlag statt Black Box.
\end{enumerate}

Der eigentliche Beitrag ist nicht die Dokumentation (bereits verbreitet), sondern die
\textbf{Prozess-Transparenz als Eingriffspunkt} sowie deren empirische Messung im Feld. Reine
Veröffentlichung von Ergebnissen liefert lediglich ein Endprodukt. Die explizite Prozessstruktur
begründet die Entscheidung und macht sie kommentierbar.

% ---------------------------------------------------------------------------
\section{Forschungsfragen}

\begin{itemize}
  \item \textbf{HF1 (Vertrauen):} Erhöht ein nachvollziehbar dargestellter Entscheidungspfad das
    Vertrauen in politische Verfahren gegenüber bloßer Ergebnisveröffentlichung?
  \item \textbf{HF2 (Transparenz):} Hilft ein expliziter Entscheidungspfad Bürgern, strittige
    oder einseitig geprägte Regelungen \emph{vor} der Abstimmung zu erkennen?
  \item \textbf{HF3 (Beteiligung, explorativ):} Führt frühe Kommentierung des Prozesses zu einer
    höheren wahrgenommenen Legitimität des Ergebnisses?
\end{itemize}

% ---------------------------------------------------------------------------
\section{Wie: Methodik}

Design Science Research mit iterativen Build-Evaluate-Zyklen, zwei Studien:
\textbf{Studie 1 (HF1)} vergleicht vier Transparenzbedingungen (B1 nur Ergebnis, B2 Endtext plus
Begründung, B3 vollständiges Protokoll, B4 interaktiver Entscheidungspfad) über verschiedene
Politikfelder; Vertrauens- und Verständnismaße werden per Befragung erhoben. \textbf{Studie 2
(HF2/HF3)} ist eine kontrollierte Studie am Interface mit gezielt verzerrten Darstellungen
(seeded bias). \textbf{Feldumgebung:} laufende Gesetzgebungs- oder Satzungsverfahren (z.\,B.\
kommunale Ebene); der Wert von Transparenz ist an abgeschlossenen, geglätteten Fällen prinzipiell
nicht messbar.

\textbf{Umsetzung und Stand (Juli 2026):} Der Aufbau des Verfahrens ist konstruktiver Bestandteil
der Arbeit. Ein lauffähiger Prototyp (Datenmodell, Prozess-Rekonstruktion, Begründungsspur
\emph{und} ein erstes interaktives Beteiligungs-Interface) belegt bereits die Machbarkeit; eine
erste Auswertung zeigt brauchbare Nachvollziehbarkeit. Im Verlauf der Arbeit werden Artefakt und
Interface zur Studienreife ausgebaut und anschließend empirisch untersucht.

\end{document}
